\begin{ex}
Trên biểu đồ tranh, mỗi biểu tượng ngôi sao tương ứng với 1 điểm tốt. Bạn Mai có 8 ngôi sao. Bạn Mai có bao nhiêu điểm tốt?
\shortans{8}
\loigiai{
Mai có 8 điểm tốt.
}
\end{ex}

\begin{ex}
Bạn Lan có 6 ngôi sao biểu thị điểm tốt. Bạn Nam có 9 ngôi sao. Cả hai bạn có bao nhiêu điểm tốt?
\shortans{15}
\loigiai{
Cả hai bạn có: $6 + 9 = 15$ điểm tốt.
}
\end{ex}

\begin{ex}
Nam nhiều hơn Lan bao nhiêu điểm tốt?
\shortans{3}
\loigiai{
Nhiều hơn: $9 - 6 = 3$ điểm tốt.
}
\end{ex}

\begin{ex}
Trên biểu đồ tranh, mỗi hình tròn biểu thị 2 cây xanh trồng được. Lớp 2A có 5 hình tròn. Lớp 2A trồng được bao nhiêu cây xanh?
\shortans{10}
\loigiai{
Lớp 2A trồng: $2 \times 5 = 10$ cây xanh.
}
\end{ex}

\begin{ex}
Lớp 2B có 7 hình tròn biểu thị cây trồng được (mỗi hình tròn là 2 cây). Lớp 2B trồng được bao nhiêu cây xanh?
\shortans{14}
\loigiai{
Lớp 2B trồng: $2 \times 7 = 14$ cây xanh.
}
\end{ex}

\begin{ex}
Lớp 2B trồng nhiều hơn lớp 2A bao nhiêu cây xanh?
\shortans{4}
\loigiai{
Nhiều hơn: $14 - 10 = 4$ cây.
}
\end{ex}

\begin{ex}
Cả hai lớp trồng được tất cả bao nhiêu cây xanh?
\shortans{24}
\loigiai{
Cả hai lớp trồng: $10 + 14 = 24$ cây.
}
\end{ex}

\begin{ex}
Biểu đồ tranh số con vật nuôi: mỗi hình con vật biểu thị 5 con. Đàn gà có 6 hình con gà. Đàn gà có bao nhiêu con gà?
\shortans{30}
\loigiai{
Có: $5 \times 6 = 30$ con gà.
}
\end{ex}

\begin{ex}
Đàn vịt có 4 hình con vịt (mỗi hình là 5 con). Đàn vịt có bao nhiêu con vịt?
\shortans{20}
\loigiai{
Có: $5 \times 4 = 20$ con vịt.
}
\end{ex}

\begin{ex}
Đàn gà nhiều hơn đàn vịt bao nhiêu con?
\shortans{10}
\loigiai{
Nhiều hơn: $30 - 20 = 10$ con.
}
\end{ex}

\begin{ex}
Cả gà và vịt có tất cả bao nhiêu con?
\shortans{50}
\loigiai{
Cả hai đàn có: $30 + 20 = 50$ con.
}
\end{ex}

\begin{ex}
Biểu đồ tranh hoa tặng cô giáo: Tổ Một có 4 bông hoa biểu tượng, Tổ Hai có 6 bông hoa biểu tượng, Tổ Ba có 5 bông hoa biểu tượng. Mỗi biểu tượng là 1 bông hoa thật. Cả 3 tổ tặng bao nhiêu bông hoa?
\shortans{15}
\loigiai{
Tặng: $4 + 6 + 5 = 15$ bông hoa.
}
\end{ex}

\begin{ex}
Tổ Hai tặng nhiều hơn Tổ Một bao nhiêu bông hoa?
\shortans{2}
\loigiai{
Nhiều hơn: $6 - 4 = 2$ bông hoa.
}
\end{ex}

\begin{ex}
Tổ nào tặng nhiều hoa nhất? Số hoa của tổ đó là bao nhiêu?
\shortans{6}
\loigiai{
Tổ Hai nhiều nhất với 6 bông hoa.
}
\end{ex}

\begin{ex}
Tổ nào tặng ít hoa nhất? Số hoa của tổ đó là bao nhiêu?
\shortans{4}
\loigiai{
Tổ Một ít nhất với 4 bông hoa.
}
\end{ex}

\begin{ex}
Biểu đồ số xe bán được trong tuần: Thứ Hai bán được 3 xe, Thứ Ba bán được 5 xe, Thứ Tư bán được 4 xe. Cả 3 ngày bán được bao nhiêu xe?
\shortans{12}
\loigiai{
Bán được: $3 + 5 + 4 = 12$ chiếc xe.
}
\end{ex}

\begin{ex}
Ngày Thứ Ba bán nhiều hơn ngày Thứ Hai bao nhiêu chiếc xe?
\shortans{2}
\loigiai{
Nhiều hơn: $5 - 3 = 2$ chiếc xe.
}
\end{ex}

\begin{ex}
Biểu đồ số bạn học sinh thích đọc sách: mỗi biểu tượng mặt cười là 2 bạn. Thể loại truyện cổ tích có 4 mặt cười. Có bao nhiêu bạn thích truyện cổ tích?
\shortans{8}
\loigiai{
Có: $2 \times 4 = 8$ bạn.
}
\end{ex}

\begin{ex}
Thể loại truyện tranh có 6 mặt cười (mỗi mặt cười là 2 bạn). Có bao nhiêu bạn thích truyện tranh?
\shortans{12}
\loigiai{
Có: $2 \times 6 = 12$ bạn.
}
\end{ex}

\begin{ex}
Số bạn thích truyện tranh nhiều hơn truyện cổ tích bao nhiêu bạn?
\shortans{4}
\loigiai{
Nhiều hơn: $12 - 8 = 4$ bạn.
}
\end{ex}

\begin{ex}
Tổng số bạn thích cả hai thể loại truyện trên là bao nhiêu bạn?
\shortans{20}
\loigiai{
Tổng số: $8 + 12 = 20$ bạn.
}
\end{ex}

\begin{ex}
Biểu đồ hoa quả bán trong ngày: Táo có 5 giỏ, Cam có 7 giỏ, Lê có 3 giỏ. Mỗi giỏ có 10 quả. Có bao nhiêu quả lê đã bán?
\shortans{30}
\loigiai{
Có: $10 \times 3 = 30$ quả lê.
}
\end{ex}

\begin{ex}
Có bao nhiêu quả cam đã bán?
\shortans{70}
\loigiai{
Có: $10 \times 7 = 70$ quả cam.
}
\end{ex}

\begin{ex}
Số quả cam bán nhiều hơn số quả táo bao nhiêu quả?
\shortans{20}
\loigiai{
Nhiều hơn: $70 - 50 = 20$ quả.
}
\end{ex}

\begin{ex}
Tổng số quả của cả 3 loại đã bán là bao nhiêu quả?
\shortans{150}
\loigiai{
Tổng số quả: $50 + 70 + 30 = 150$ quả.
}
\end{ex}
